\documentclass[11pt]{article}
\usepackage[a4paper,margin=1in]{geometry}
\usepackage{amsmath,amssymb,amsthm,booktabs}
\usepackage[T1]{fontenc}
\usepackage{lmodern}
\usepackage[hidelinks]{hyperref}
\setlength{\emergencystretch}{3em}

\theoremstyle{plain}
\newtheorem{theorem}{Theorem}
\newtheorem{lemma}[theorem]{Lemma}
\newtheorem{proposition}[theorem]{Proposition}
\newtheorem{corollary}[theorem]{Corollary}
\theoremstyle{remark}
\newtheorem{remark}[theorem]{Remark}

\newcommand{\dcut}{\delta_{\square}}

\title{A Counterexample to Conjecture 00000000046}
\author{%
  Feiyu\thanks{Independent researcher.}
  \and
  Claude (Opus 5)\thanks{Anthropic.}%
}
\date{2026-09-12}

\begin{document}
\maketitle

\begin{abstract}
Conjecture 00000000046 asserts that the prime-sum graphs $G_n$ converge, in cut
distance, to the constant graphon $W\equiv1/2$. They do not, and not merely
slowly: every $G_n$ is triangle-free, so its triangle density is $0$ while that
of $W$ is $1/8$, and the counting lemma puts $\dcut(W_{G_n},W)\ge1/24$ for
\emph{every} $n$. The edge density fails too, tending to $0$ rather than $1/2$.
The sequence does converge --- to the zero graphon.
\end{abstract}

\section{Statement}

For $n\ge1$ let $G_n$ be the graph on $\{1,\ldots,n\}$ with $i\sim j$ if and
only if $i\ne j$ and $i+j$ is prime. Let $W_{G_n}$ denote the step graphon
associated with $G_n$, let $\dcut$ be the cut distance, and for a finite simple
graph $F$ let $t(F,\cdot)$ be homomorphism density. Conjecture 00000000046
asserts
\[
  \dcut\bigl(W_{G_n},\,\tfrac12\bigr)\longrightarrow0 ,
\]
where $\tfrac12$ denotes the constant graphon $W\equiv1/2$.

\section{\texorpdfstring{$G_n$}{Gn} is bipartite, hence triangle-free}

\begin{lemma}\label{lem:bip}
Let $i\ne j$ be positive integers of the same parity. Then $i\sim j$ fails in
every $G_n$.
\end{lemma}

\begin{proof}
$i+j$ is even. Being distinct positive integers, $i+j\ge1+3=4$ when both are
odd and $i+j\ge2+4=6$ when both are even; in either case $i+j>2$. An even
number greater than $2$ is not prime.
\end{proof}

So every edge of $G_n$ joins an odd vertex to an even one: $G_n$ is bipartite
with sides the odd and the even vertices.

\begin{theorem}\label{thm:tri}
$G_n$ is triangle-free for every $n$.
\end{theorem}

\begin{proof}
Among any three vertices, two have the same parity; by Lemma~\ref{lem:bip}
they are not adjacent.
\end{proof}

\section{The conjecture is false, quantitatively}

\begin{theorem}\label{thm:main}
For every $n$,
\[
  \dcut\bigl(W_{G_n},\tfrac12\bigr)\;\ge\;\frac1{24} .
\]
In particular $\dcut(W_{G_n},\tfrac12)\not\to0$ and Conjecture 00000000046 is
false.
\end{theorem}

\begin{proof}
Since $G_n$ has no loops, every homomorphism $K_3\to G_n$ is injective, so
$t(K_3,W_{G_n})=6\cdot\#\{\text{triangles of }G_n\}/n^3$, which is $0$ by
Theorem~\ref{thm:tri}. For the constant graphon,
$t(K_3,\tfrac12)=(\tfrac12)^3=\tfrac18$. The counting lemma states
\[
  \bigl|t(F,U)-t(F,U')\bigr|\;\le\;e(F)\,\dcut(U,U')
\]
for graphons $U,U'$ and finite simple $F$. Taking $F=K_3$, so $e(F)=3$,
\[
  \tfrac18=\bigl|t(K_3,\tfrac12)-t(K_3,W_{G_n})\bigr|\le3\,\dcut\bigl(W_{G_n},\tfrac12\bigr),
\]
which is the claim.
\end{proof}

\section{The edge density fails as well}

Triangle-freeness already settles the conjecture. The first moment fails too,
and in the opposite direction from what $W\equiv1/2$ would require.

\begin{proposition}\label{prop:dens}
Let $d(G_n)=e(G_n)\big/\binom n2$. Then
\[
  d(G_n)\;\le\;\frac{n+5}{3(n-1)}+\frac{2}{n(n-1)}\;=\;\frac13+O\!\left(\frac1n\right).
\]
\end{proposition}

\begin{proof}
Let $i<j\le n$ with $i+j$ prime. By Lemma~\ref{lem:bip}, $i+j$ is odd. If
$3\mid i+j$ then $i+j$, being prime, equals $3$, which happens only for the
single pair $\{1,2\}$. Every other edge therefore has
$i+j\equiv1$ or $5\pmod6$. For fixed $i$ and a fixed residue $r$, the number of
$j\le n$ with $i+j\equiv r\pmod 6$ is at most $\lceil n/6\rceil\le(n+5)/6$; two
residues give at most $(n+5)/3$. Summing over $i\le n$ counts each edge twice,
so $e(G_n)\le n(n+5)/6+1$, and dividing by $\binom n2=n(n-1)/2$ gives the
bound.
\end{proof}

In fact $d(G_n)\to0$: each prime $p\le2n$ arises as $i+j$ for at most $n/2$
pairs $i<j\le n$, so $e(G_n)\le n\pi(2n)/2$ and
$d(G_n)\le n\pi(2n)/\bigl(n(n-1)\bigr)=O(1/\log n)$ by the prime number
theorem. The computed values agree:

\begin{center}
\begin{tabular}{lrrrrrrrr}
\toprule
$n$ & 10 & 20 & 50 & 100 & 200 & 400 & 800 & 1600\\
\midrule
$e(G_n)$ & 18 & 56 & 295 & 1044 & 3676 & 13239 & 47811 & 174317\\
$d(G_n)$ & .400 & .295 & .241 & .211 & .185 & .166 & .150 & .136\\
\bottomrule
\end{tabular}
\end{center}

\begin{remark}[what the sequence does converge to]
Since $W_{G_n}\ge0$, the cut norm is at most the $L^1$ norm, so
\[
  \dcut\bigl(W_{G_n},0\bigr)\le\int W_{G_n}=d(G_n)+O(1/n)\longrightarrow0 .
\]
Thus $G_n$ converges in cut distance to the \emph{zero} graphon $W\equiv0$, not
to $W\equiv1/2$. A sequence of graphs converges to a constant graphon $c$ only
if its edge density tends to $c$; here that forces $c=0$. The conjecture would
be correct only for a sequence of density $1/2$, which $G_n$ is not, and the
triangle obstruction of Theorem~\ref{thm:main} rules out $W\equiv1/2$ on its
own.
\end{remark}

\section{Formal verification}

The accompanying Lean~4 project formalizes the combinatorial core against
Mathlib:

\begin{itemize}
  \item \texttt{Adj} --- adjacency in the prime-sum graph, $i\ne j$ and
    $i+j$ prime, with primality Mathlib's \texttt{Nat.Prime}.
  \item \texttt{not\_adj\_of\_same\_parity}, \texttt{parity\_ne\_of\_adj} ---
    Lemma~\ref{lem:bip}: same-parity vertices are never adjacent, so $G_n$ is
    bipartite between odds and evens.
  \item \texttt{no\_triangle} --- Theorem~\ref{thm:tri}: no three vertices are
    pairwise adjacent.
  \item \texttt{triangleCount}, \texttt{triangleCount\_eq\_zero} --- the
    triangle count of $G_n$, over increasing triples of $\{1,\ldots,n\}$, is
    $0$ for every $n$.
\end{itemize}

\paragraph{Scope.} Mathlib contains no graph-limit theory: there are no
graphons, no cut norm or cut distance, and no homomorphism densities. The step
from Theorem~\ref{thm:tri} to Theorem~\ref{thm:main} --- the counting lemma
--- is therefore carried out in this paper and is \emph{not} part of the Lean
development, which establishes the input to it. We state this explicitly rather
than leave the coverage of the formalization to be inferred. Building the
missing theory is a substantial project in its own right and is not attempted
here.

The project builds with \texttt{lake build} on
\texttt{leanprover/lean4:v4.33.1} with Mathlib \texttt{v4.33.1}. It contains no
\texttt{sorry}, no \texttt{admit} and no \texttt{native\_decide}, and
introduces no axioms of its own:
\texttt{\#print axioms triangleCount\_eq\_zero} reports exactly the three
standard axioms \texttt{propext}, \texttt{Classical.choice} and
\texttt{Quot.sound}.

\end{document}
